% !TeX encoding = UTF-8
\documentclass[variant=guia,cover=institutional,mode=screen]{ua-engineering-v6-article}
% !TeX encoding = UTF-8
\UASetUniversity{Universidad Austral}
\UASetFaculty{Facultad de Ingeniería}
\UASetDepartment{Departamento de Computación}
\UASetCampus{Pilar}
\UASetCareer{Ingeniería Informática}
\UASetCommission{Comisión A}
\UASetProfessor{Equipo docente}
\UASetSemester{Segundo cuatrimestre 2026}
\UASetCourse{Sistemas Distribuidos}
\UASetDate{2026}


\UASetClassNumber{13}
\UASetClassTitle{Confiabilidad operativa}
\UASetDocKind{Guía de ejercicios}

\title{Clase 13 --- Guía de ejercicios}
\author{\UAProfessor}
\date{\UADate}

\begin{document}

\section{Indicaciones generales}
Resuelva cada ejercicio declarando supuestos, modelo de fallas, garantía y trade-off. Cuando corresponda, construya una ejecución verificable y vincule la respuesta con evidencia observable.

\section{Nivel 1: fundamentos y reconocimiento}
\begin{uaexercise}
Defina SLI, SLO y SLA, y explique la diferencia entre los tres conceptos.
\end{uaexercise}

\begin{uaexercise}
Explique por qué un SLI debe elegirse desde la perspectiva del usuario y no solo desde la infraestructura.
\end{uaexercise}

\begin{uaexercise}
Dé dos ejemplos de buenos SLI y dos ejemplos de malos SLI.
\end{uaexercise}

\begin{uaexercise}
Explique por qué “CPU media del cluster” no siempre es un buen SLI.
\end{uaexercise}

\begin{uaexercise}
Diseñe un SLO razonable para una API pública de lectura.
\end{uaexercise}


\section{Nivel 2: ejecuciones y fallas}
\begin{uaexercise}
Diseñe un SLO razonable para un sistema de checkout.
\end{uaexercise}

\begin{uaexercise}
Explique por qué un SLO de 100\% suele ser una mala idea práctica.
\end{uaexercise}

\begin{uaexercise}
Explique por qué elegir una ventana temporal adecuada es parte del diseño del SLO.
\end{uaexercise}

\begin{uaexercise}
Defina error budget y explique su significado operativo.
\end{uaexercise}

\begin{uaexercise}
Calcule el error budget mensual para un servicio con SLO de 99.9\%.
\end{uaexercise}


\section{Nivel 3: diseño y comparación}
\begin{uaexercise}
Explique qué decisiones podría tomar un equipo si está consumiendo demasiado rápido su error budget.
\end{uaexercise}

\begin{uaexercise}
Explique por qué el error budget conecta confiabilidad con velocidad de cambio.
\end{uaexercise}

\begin{uaexercise}
Explique por qué un SLO basado en promedio de latencia puede ser engañoso.
\end{uaexercise}

\begin{uaexercise}
Diseñe un SLI/SLO usando p95 o p99 para un servicio de pagos.
\end{uaexercise}

\begin{uaexercise}
Explique qué problema aparece si el equipo mide algo distinto de lo que el usuario experimenta.
\end{uaexercise}


\section{Nivel 4: integración y defensa}
\begin{uaexercise}
Diseñe una estrategia mínima para medir correctamente un SLI de disponibilidad.
\end{uaexercise}

\begin{uaexercise}
Explique qué significa operar un sistema guiado por SLOs.
\end{uaexercise}

\begin{uaexercise}
Construya un ejemplo donde un incidente técnico parezca grave, pero no implique violación del SLO.
\end{uaexercise}

\begin{uaexercise}
Construya un ejemplo donde un problema “pequeño” técnicamente sí implique un daño real al servicio medido por SLO.
\end{uaexercise}

\begin{uaexercise}
Explique cómo SLOs ayudan a resolver discusiones entre producto e ingeniería.
\end{uaexercise}


\end{document}
